% =====================================================================
%  BAB III METODOLOGI
% =====================================================================

\chapter{METODOLOGI}

\section{Rancangan Penelitian}

Penelitian ini merupakan penelitian rekayasa yang membangun sistem pembacaan
tampilan tujuh segmen dengan lapisan verifikasi kontekstual. Alur kerja dibagi
menjadi lima tahap, yaitu:
\begin{urut}
  \item penetapan spesifikasi data dan metrik keberhasilan;
  \item pembangkitan data latih sintetis dan pengumpulan data uji nyata;
  \item pengukuran kemampuan awal PP-OCRv6 tanpa adaptasi;
  \item pembangunan lapisan verifikasi secara bertahap dan terukur;
  \item evaluasi sistem penuh serta analisis dampaknya terhadap operator.
\end{urut}

Setiap lapisan verifikasi dibangun dan diukur secara terpisah. Pemisahan ini
disengaja agar penyebab perubahan laju kesalahan senyap dapat diidentifikasi,
sesuatu yang tidak mungkin dilakukan bila beberapa lapisan diuji sekaligus.

\section{Data dan Instrumen}

\subsection{Data Latih Sintetis}

Data latih dibangkitkan secara terprogram dengan variasi terkontrol, meliputi
jenis font segmen, warna dan intensitas pancaran, ketebalan segmen, kontras
tampilan \textit{liquid crystal display} pasif, sudut kamera, derajat
kemiringan, derau sensor, serta panjang bacaan 1 sampai 8 karakter. Prinsip yang
dipedomani adalah variasi lebih penting daripada jumlah; citra yang banyak namun
seragam tidak menambah kemampuan generalisasi. Pemanfaatan pembangkit sintetis
untuk keperluan serupa telah dilakukan \textcite{gushima2023automatic} dan
\textcite{yenneti2023lightweight}.

Rincian jumlah citra, sebaran kondisi, dan format label ditetapkan pada dokumen
spesifikasi data tersendiri. Data latih tidak pernah dipakai sebagai data
penilaian.

\subsection{Data Uji Nyata}

Data uji berupa citra timbangan digital nyata yang diambil dari sejumlah merek
berbeda pada beberapa kondisi pencahayaan dan sudut pengambilan. Setiap
pembacaan disertai rekaman video pendek agar lapisan verifikasi temporal dapat
diuji. Kumpulan data uji dikunci sebelum eksperimen dijalankan, dan setelah
dikunci tidak boleh dipakai untuk memilih ambang, menyetel praproses, maupun
menentukan hyperparameter.

Ukuran kumpulan data uji dan konsekuensinya terhadap lebar interval kepercayaan
dilaporkan secara eksplisit. Pengumpulan data nyata tetap diperlukan karena
tidak tersedia kumpulan data publik berupa timbangan dengan label tingkat string
beserta titik desimal dan metadata kondisi.

\subsection{Instrumen}

Perangkat lunak yang dipakai adalah PP-OCRv6 \parencite{zhang2026ppocrv6} yang
dijalankan melalui ONNX Runtime agar dapat beroperasi tanpa koneksi jaringan,
serta pustaka pengolahan citra untuk praproses. Sebagai pembanding dipakai
CRNN-CTC \parencite{shi2017crnn} dan perangkat klasik \textit{seven segment
optical character recognition} \parencite{auerswal_ssocr}.

\section{Rancangan Lapisan Verifikasi}

Sistem menerima citra, memotong area tampilan, melakukan praproses, lalu
menjalankan pengenalan. Keluaran pengenalan tidak langsung diterima, melainkan
melewati empat lapisan verifikasi seperti tampak pada Gambar~\ref{fig:alur}.

\begin{figure}[H]
\centering
\begin{tikzpicture}[
  node distance=0.5cm,
  kotak/.style={rectangle, draw, rounded corners=2pt, align=center,
                text width=9.6cm, inner sep=4pt, font=\small},
  lapis/.style={rectangle, draw, thick, rounded corners=2pt, align=center,
                text width=9.6cm, inner sep=4pt, font=\small},
  akhir/.style={rectangle, draw, very thick, rounded corners=2pt,
                align=center, text width=3.6cm, inner sep=5pt, font=\small},
  panah/.style={-{Latex[length=2mm]}, thick}
]
\node[kotak] (aku) {Akuisisi citra atau 3--5 bingkai video};
\node[kotak, below=of aku] (pra) {Pemotongan area tampilan dan praproses};
\node[kotak, below=of pra] (ocr) {PP-OCRv6: deteksi dan pengenalan};
\node[kotak, below=of ocr] (kan) {Kandidat string dan nilai keyakinan};

\node[lapis, below=0.7cm of kan] (l1)
  {\textbf{Lapis 1 --- Verifikasi Struktural}\\[1pt]
   {\scriptsize gagal $\rightarrow$ tolak}};
\node[lapis, below=of l1] (l2)
  {\textbf{Lapis 2 --- Verifikasi Fisik}\\[1pt]
   {\scriptsize gagal $\rightarrow$ tolak}};
\node[lapis, below=of l2] (l3)
  {\textbf{Lapis 3 --- Verifikasi Temporal}\\[1pt]
   {\scriptsize bingkai tidak sepakat $\rightarrow$ minta konfirmasi}};
\node[lapis, below=of l3] (l4)
  {\textbf{Lapis 4 --- Verifikasi Statistik}\\[1pt]
   {\scriptsize keyakinan rendah $\rightarrow$ minta konfirmasi}};

\node[akhir, below=0.7cm of l4] (terima) {\textbf{TERIMA}};

\draw[panah] (aku) -- (pra);
\draw[panah] (pra) -- (ocr);
\draw[panah] (ocr) -- (kan);
\draw[panah] (kan) -- (l1);
\draw[panah] (l1) -- node[right, font=\scriptsize, pos=0.45] {lolos} (l2);
\draw[panah] (l2) -- node[right, font=\scriptsize, pos=0.45] {lolos} (l3);
\draw[panah] (l3) -- node[right, font=\scriptsize, pos=0.45] {lolos} (l4);
\draw[panah] (l4) -- node[right, font=\scriptsize, pos=0.45] {lolos} (terima);
\end{tikzpicture}
\caption{Alur pembacaan dengan empat lapisan verifikasi kontekstual}
\label{fig:alur}
\end{figure}

\subsection{Lapis 1: Verifikasi Struktural}

Lapisan ini memeriksa kesesuaian jumlah digit yang terbaca dengan jumlah slot
yang tersedia pada tampilan, serta kesahihan posisi titik desimal. Pemeriksaan
ini menyasar pola kegagalan yang bersifat struktural, yaitu hilangnya karakter
berulang pada dekoder CTC \parencite{graves2006ctc,oicws2026lcd}.

\subsection{Lapis 2: Verifikasi Fisik}

Lapisan ini memeriksa kesesuaian bacaan dengan konteks domain, meliputi rentang
berat yang sah untuk suatu produk, toleransi penimbangan, dan selisih terhadap
riwayat transaksi sebelumnya. Konfigurasi produk disimpan pada basis data dan
menjadi rujukan keputusan.

\subsection{Lapis 3: Verifikasi Temporal}

Lapisan ini membandingkan hasil pembacaan pada beberapa bingkai dari pembacaan
yang sama. Bacaan yang benar cenderung stabil, sedangkan kesalahan cenderung
berubah-ubah. Gagasan penyaringan antarbingkai telah dipakai untuk stabilisasi
keluaran \parencite{yenneti2023lightweight}, tetapi belum dimanfaatkan sebagai
dasar keputusan menerima atau menolak.

\subsection{Lapis 4: Verifikasi Statistik}

Lapisan ini memakai nilai keyakinan yang telah dikalibrasi. Kalibrasi dilakukan
pada bagian data yang terpisah dari data penilaian, sehingga ambang keputusan
tidak ditentukan dari data yang dipakai menilai kinerja.

\section{Rancangan Sistem}

Keputusan setiap lapisan digabungkan menjadi tiga kemungkinan keluaran, yaitu
terima, tolak, atau minta konfirmasi. Keluaran terima dan konfirmasi disimpan ke
basis data transaksi beserta alasannya, sedangkan keluaran tolak dicatat pada log
penolakan. Semua keputusan meninggalkan rekaman sehingga dapat diaudit dan
menjadi bahan perhitungan laju kesalahan senyap serta tingkat koreksi operator.
Sistem dirancang berjalan tanpa koneksi jaringan dan melakukan sinkronisasi
ketika jaringan tersedia.

\section{Rancangan Eksperimen}

Eksperimen disusun berjenjang agar kontribusi setiap tahap dapat dipisahkan.
Rincian disajikan pada Tabel~\ref{tab:eksperimen}.

\begin{table}[H]
\caption{Rancangan eksperimen}
\label{tab:eksperimen}
\begin{tabularx}{\textwidth}{@{}p{0.8cm} p{5.2cm} X@{}}
\toprule
\textbf{Kode} & \textbf{Konfigurasi} & \textbf{Pertanyaan yang dijawab} \\
\midrule
E1 & PP-OCRv6 \textit{tiny}, \textit{small}, dan \textit{medium} tanpa adaptasi & Sejauh mana kemampuan tanpa adaptasi? \\
E2 & Penggantian kepala keluaran menjadi 12 kelas & Tingkat adaptasi minimum yang memadai? \\
E3 & Ditambah verifikasi struktural & Berapa kesalahan senyap yang tertangkap? \\
E4 & Ditambah verifikasi fisik & Tambahan manfaat dari konteks produk? \\
E5 & Ditambah verifikasi temporal dan kalibrasi & Di mana titik seimbang antara \textit{coverage} dan akurasi? \\
E6 & Sistem penuh dibandingkan kondisi awal & Dampak terhadap tingkat koreksi dan waktu per transaksi \\
\bottomrule
\end{tabularx}
\end{table}

\section{Metrik Evaluasi}

Metrik yang dipakai diuraikan sebagai berikut.
\begin{poin}
  \item \textit{Full-string recognition} (FSR), yaitu persentase bacaan yang
        benar secara utuh.
  \item \textit{Character error rate} (CER), yaitu laju kesalahan tingkat
        karakter.
  \item Akurasi posisi desimal, yaitu persentase titik desimal pada posisi yang
        benar.
  \item Laju kesalahan senyap, yaitu persentase bacaan salah yang diterima
        sistem tanpa peringatan. Metrik ini merupakan metrik utama.
  \item \textit{Coverage}, yaitu persentase bacaan yang diterima otomatis.
  \item \textit{Selective accuracy}, yaitu akurasi pada bacaan yang diterima.
  \item AURC, yaitu luas daerah di bawah kurva \textit{risk--coverage}.
  \item Latensi p50 dan p95, ukuran model, serta tingkat koreksi operator.
\end{poin}

Setiap persentase yang dilaporkan disertai jumlah sampel dan interval
kepercayaan 95 persen. Perbandingan latensi disertai keterangan ukuran model dan
perangkat yang dipakai.

\section{Jadwal Penelitian}

Pelaksanaan penelitian direncanakan dalam sepuluh minggu kerja sebagai berikut.
\begin{poin}
  \item Minggu 1: penetapan spesifikasi data, metrik, dan ambang keberhasilan.
  \item Minggu 2: pembangunan pembangkit data sintetis dan penyusunan data uji
        terkontrol.
  \item Minggu 3: pengukuran kemampuan awal PP-OCRv6 untuk eksperimen E1.
  \item Minggu 4: pengambilan data uji nyata, penguncian data uji, dan
        pelaksanaan E2.
  \item Minggu 5: implementasi dan pengukuran verifikasi struktural (E3).
  \item Minggu 6: implementasi dan pengukuran verifikasi fisik (E4).
  \item Minggu 7: verifikasi temporal, kalibrasi, dan penyusunan kurva
        \textit{risk--coverage} (E5).
  \item Minggu 8: perangkaian sistem penuh dan pengujian fungsional (E6).
  \item Minggu 9: analisis kesalahan dan penyusunan Bab IV.
  \item Minggu 10: penyusunan Bab V dan konsolidasi naskah.
\end{poin}
